\documentclass[11pt]{article}
\usepackage[a4paper,margin=20mm]{geometry}
\usepackage{fontspec}
\usepackage{xeCJK}
\usepackage{array}
\usepackage{booktabs}
\usepackage{longtable}
\usepackage{tabularx}
\usepackage{enumitem}
\usepackage{xcolor}
\usepackage{hyperref}
\usepackage{setspace}

\setmainfont{PingFang SC}
\setsansfont{PingFang SC}
\setCJKmainfont{PingFang SC}
\setCJKsansfont{PingFang SC}
\setmonofont{Menlo}
\setlength{\parskip}{6pt}
\setlength{\parindent}{0pt}
\renewcommand{\arraystretch}{1.2}
\setlist[itemize]{leftmargin=18pt,itemsep=3pt,topsep=3pt}
\hypersetup{colorlinks=true,linkcolor=blue,urlcolor=blue}

\definecolor{brandblue}{HTML}{2F6BFF}
\definecolor{textgray}{HTML}{4A5568}

\begin{document}

{\LARGE\bfseries benchmark\_v2 有效性审查与复跑报告}\\
{\small \color{textgray} 日期：2026-05-05 \quad 目标：说明 benchmark 当前到底测到了什么、没测到什么，并给出一轮新的本地结果}

\section{报告定位}

这份报告不是系统总验收，也不是单纯展示一个通过率。

它分四层叙事：
\begin{itemize}
\item 结构层：fixture / runner / reporting 是否可用
\item 轻量回归层：不依赖本地 LLM 写入链时，benchmark 能否稳定暴露实时分类和干扰误判
\item 深度子集层：开启 \texttt{FULL\_WRITE + DEEP\_EVAL} 后，写入侧卡片检查是否能打出真实失败
\item 有效性层：按公开 memory benchmark 常见原则，判断这些结果可以支持哪些结论，不能支持哪些结论
\end{itemize}

\section{审查依据}

本次有效性审查参考了近年的几类 memory benchmark 设计原则：
\begin{itemize}
\item LongMemEval：把长期记忆拆成 extraction、multi-session reasoning、temporal reasoning、updates、abstention 等能力，而不是只看单一命中率
\item EverMemBench：强调 multi-party、multi-group、长时间跨度、跨主题交错和角色差异
\item MemBench：强调不要只看 accuracy，还要看效率、容量和使用场景差异
\item StructMemEval：提醒不能只测“记没记住”，还要看记忆结构是否组织得合理
\item MemoryAgentBench：强调 incremental multi-turn interaction 和 selective forgetting / update
\end{itemize}

因此，本报告重点检查四件事：
\begin{itemize}
\item 是否存在负样本、近似误判样本、跨群漂移样本
\item 是否区分轻量路由能力与深度写入能力
\item 是否把 long-term recall、update、evidence、relation 分层看待
\item 是否避免把 filtered run 或轻量代理指标误写成总能力结论
\end{itemize}

\section{本轮运行内容}

\small
\begin{tabularx}{\textwidth}{@{}>{\raggedright\arraybackslash}p{1.9cm}XX@{}}
\toprule
\textbf{层次} & \textbf{作用} & \textbf{本轮输出} \\
\midrule
结构校验 & 检查 source/runtime schema 是否一致 & \texttt{validate\_fixture} 通过 \\
单元测试 & 检查 runner、reporting、replay adapter 兼容性 & 11/11 通过 \\
轻量全量回放 & 主看实时分类、忽略规则、困难样本区分能力 & 生成 \texttt{light\_full\_results.json} \\
深度子集回放 & 主看 FULL\_WRITE 后的卡片命中与环境退化信号 & 生成 \texttt{deep\_oc\_launch\_ops\_results.json} \\
\bottomrule
\end{tabularx}
\normalsize

本轮实际执行命令：
\begingroup
\small\ttfamily
conda run -n feishu-ai-p0 python benchmark\_v2/validate\_fixture.py

conda run -n feishu-ai-p0 python -m unittest \\
tests.test\_benchmark\_v2\_runner tests.test\_benchmark\_replay

conda run -n feishu-ai-p0 python -m benchmark\_v2.dual\_channel\_runner \\
benchmark\_v2/full\_demo\_case\_v2.json

BENCHMARK\_V2\_DEEP\_EVAL=1 FULL\_WRITE=1 \\
conda run -n feishu-ai-p0 python -m benchmark\_v2.dual\_channel\_runner \\
benchmark\_v2/full\_demo\_case\_v2.json --chat oc\_launch\_ops
\par
\endgroup

\section{当前样本池分类}

当前 runtime fixture 共 24 个 batch。下面的分类是\textbf{标签覆盖计数}，不是互斥分桶，所以相加会超过 24。

\begin{tabularx}{\textwidth}{@{}l>{\raggedleft\arraybackslash}p{2cm}X@{}}
\toprule
\textbf{类别} & \textbf{覆盖 batch} & \textbf{说明} \\
\midrule
结构 / 分块样本 & 2 & progress、sparse window / lookback \\
干扰 / 路由样本 & 15 & noise、query、schedule、task、parallel、topic\_boundary、near\_miss 等 \\
冲突 / 更新样本 & 4 & refine、supersede、conflict、long-term update \\
长线追问样本 & 7 & source、version、summary、topic\_list、long-term recall \\
非工程协作样本 & 5 & 发布运营、客服口径、会务组织、答辩物料 \\
\bottomrule
\end{tabularx}

本轮新增后的困难样本特征：
\begin{itemize}
\item 带时间词但本质是发布策略或风险口径，而不是 schedule
\item 非 \texttt{@机器人} 的人对人复述请求，防止 query 过敏
\item 直播备用链接预案，防止误判成日程
\item 显式排除后续改口的 long-term recall
\item 同一时间窗口内碎片化多线程收口，防止 topic merge 过度
\end{itemize}

\section{当前会生成哪些指标}

\subsection{结构层}

\begin{tabularx}{\textwidth}{@{}lX@{}}
\toprule
\textbf{指标} & \textbf{意义} \\
\midrule
\texttt{validate\_fixture} 是否通过 & source/runtime schema、source-only 字段剥离、基本字段完整性 \\
单元测试通过数 & runner / reporting / replay adapter 的局部可用性 \\
\bottomrule
\end{tabularx}

\subsection{轻量全量层}

\begin{tabularx}{\textwidth}{@{}lX@{}}
\toprule
\textbf{指标} & \textbf{意义} \\
\midrule
\texttt{passed\_batches / total\_batches} & 当前样本包上的总体通过率 \\
\texttt{interference\_metrics.batch\_pass\_rate} & 干扰相关 batch 总体通过率 \\
\texttt{realtime\_action\_match\_rate} & 实时动作是否符合 fixture 标注 \\
\texttt{write\_count\_match\_rate} & 写入侧产出数量是否符合轻量预期 \\
\texttt{ignore\_rule\_match\_rate} & 被标记为应忽略的消息是否确实被隔离 \\
\texttt{difficult / multi\_intent / near\_miss batch pass rate} & 困难样本、混合意图样本、近似误判样本的区分能力 \\
\texttt{realtime\_action\_distribution} & 当前 fixture 下动作分布是否失衡 \\
\texttt{performance.*} & 回放延迟、写入延迟、吞吐 \\
\bottomrule
\end{tabularx}

\subsection{深度子集层}

\begin{tabularx}{\textwidth}{@{}lX@{}}
\toprule
\textbf{指标} & \textbf{意义} \\
\midrule
\texttt{expected\_memory\_cards[*]} & 写入后卡片内容是否命中预期关键词 \\
\texttt{expected\_relations[*]} & 关系抽取是否命中预期类型和前后卡片 \\
\texttt{expected\_evidence\_checks[*]} & 来源 / 证据是否能回到预期 message ids \\
\texttt{avg / p95 write latency} & FULL\_WRITE 真实耗时 \\
Graphiti / Stage 2 退化日志 & 区分环境降级和系统能力失败 \\
\bottomrule
\end{tabularx}

\section{有效性审查}

\subsection{1. Construct Validity：当前 benchmark 是否在测对的东西}

结论：\textbf{基本成立，但只对一部分能力成立。}

成立的部分：
\begin{itemize}
\item 当前 fixture 已经能测到实时分类的现实误判，不再只是功能 smoke test
\item 已经有 near-miss、cross-group、long-gap recall、non-engineering 场景
\item 轻量层和深度层的能力口径已经分开
\end{itemize}

仍然不足的部分：
\begin{itemize}
\item conflict / relation 目前在轻量全量模式下只统计 batch 覆盖，不统计真实 relation 命中
\item case 级 final recall 这轮没有跑完整 deep，因此不能宣称终态 long-term retrieval 已验证
\item 当前 fixture 仍是单语言、单主项目家族，外部有效性有限
\end{itemize}

\subsection{2. Discriminative Validity：当前 benchmark 能不能把强弱区分开}

结论：\textbf{可以，但目前主要体现在 realtime routing，而不是深度 relation。}

证据是这轮轻量全量结果里：
\begin{itemize}
\item 总体通过率不是 100\%，而是 21/24
\item 干扰相关 batch 通过率是 13/16 = 0.8125
\item multi-intent 样本通过率是 5/7 = 0.7143
\item near-miss 样本通过率是 3/6 = 0.5
\end{itemize}

这说明样本已经能把“普通规则命中”和“现实近似误判”拉开。

\subsection{3. Internal Validity：结果是否会被口径混淆}

结论：\textbf{比之前清楚了，但仍需在报告里明确限制条件。}

需要明确写清楚的点：
\begin{itemize}
\item 标签分类是重叠计数，不是互斥分桶
\item filtered deep run 会跳过 case-level eval，不能拿来证明完整终态 recall
\item \texttt{deep\_eval\_disabled} 时，relation 命中率应视为 skipped，而不是 0
\item Graphiti 未初始化会让 Stage 2 降级到 Jaccard，这种结果不能和完整环境混读
\end{itemize}

\subsection{4. External Validity：离真实生产还有多远}

结论：\textbf{离“生产级评测体系”更近了，但还不是完整生产外推。}

当前优点：
\begin{itemize}
\item 已包含多群、多角色、非工程协作、跨群漂移和长线追问
\item 已显式引入本体构造、时间基底、线程槽位和 iceberg 隐层
\end{itemize}

当前不足：
\begin{itemize}
\item 样本规模仍小，只有 24 个 batch
\item 中文单域为主，缺少多语言和多租户
\item FULL\_WRITE 深度评测还没有稳定跑完整 case
\item relation / evidence / topic summary 的深度成功率还没有形成完整统计面
\end{itemize}

\section{复跑结果}

\subsection{A. 结构层结果}

\begin{tabularx}{\textwidth}{@{}l>{\raggedleft\arraybackslash}p{2cm}X@{}}
\toprule
\textbf{检查项} & \textbf{结果} & \textbf{说明} \\
\midrule
fixture 校验 & 通过 & source/runtime schema 一致，source-only 字段剥离正常 \\
单元测试 & 11/11 & runner、reporting、replay adapter 基本可用 \\
\bottomrule
\end{tabularx}

\subsection{B. 轻量全量回放结果}

\begin{tabularx}{\textwidth}{@{}l>{\raggedleft\arraybackslash}p{2.4cm}>{\raggedleft\arraybackslash}p{2.4cm}X@{}}
\toprule
\textbf{指标} & \textbf{结果} & \textbf{分母} & \textbf{解释} \\
\midrule
总体 batch 通过率 & 87.5\% & 21/24 & 不再是全绿，说明样本已有区分度 \\
failed checks & 3 & 70 项检查 & 失败集中在 realtime routing \\
干扰相关 batch 通过率 & 81.25\% & 13/16 & 干扰样本已能稳定暴露问题 \\
realtime action 匹配率 & 81.25\% & 13/16 & 失败点就是当前规则分类缺口 \\
write count 匹配率 & 100\% & 16/16 & 轻量写入数量口径基本稳定 \\
ignore rule 命中率 & 100\% & 5/5 & 负样本隔离目前稳定 \\
difficult batch 通过率 & 100\% & 6/6 & 当前“难”的定义仍偏结构难，不完全等于语义难 \\
multi-intent batch 通过率 & 71.43\% & 5/7 & 混合意图仍明显拉低成功率 \\
near-miss batch 通过率 & 50.0\% & 3/6 & 当前最有价值的区分信号 \\
\bottomrule
\end{tabularx}

性能口径：
\begin{itemize}
\item case runtime: 6.559 ms
\item avg realtime latency: 0.018 ms
\item avg write latency: 0.080 ms
\item p95 write latency: 0.137 ms
\end{itemize}

这层能得出的结论：
\begin{itemize}
\item benchmark 已经能稳定暴露实时分类误判
\item 误判集中在“带时间词的发布/风险口径”与“风险预案”
\item 这层\textbf{不能}证明 relation、evidence、终态 recall 成功
\end{itemize}

\subsection{轻量层失败 batch}

\begin{tabularx}{\textwidth}{@{}lX@{}}
\toprule
\textbf{batch} & \textbf{失败含义} \\
\midrule
\texttt{batch\_016\_product\_launch\_ops} & “晚上八点灰度开放给老师试用”被规则判成 \texttt{schedule}，但 benchmark 认为它应属于发布策略口径 \\
\texttt{batch\_017\_meeting\_org\_external} & “直播挂了就切备用腾讯会议链接”被规则判成 \texttt{schedule}，但 benchmark 认为它是风险预案 / 会务规则 \\
\texttt{batch\_022\_release\_policy\_time\_boundary} & 同一批中“晚上八点灰度开放”与“今晚七点前补完 FAQ”被判成 \texttt{schedule + schedule}，而 benchmark 期望 \texttt{policy + task} \\
\bottomrule
\end{tabularx}

\subsection{C. 深度子集回放结果（\texttt{oc\_launch\_ops}）}

\begin{tabularx}{\textwidth}{@{}l>{\raggedleft\arraybackslash}p{2.4cm}>{\raggedleft\arraybackslash}p{2.4cm}X@{}}
\toprule
\textbf{指标} & \textbf{结果} & \textbf{分母} & \textbf{解释} \\
\midrule
batch 通过率 & 0\% & 0/3 & 三个发布运营子样本全部暴露问题 \\
failed checks & 5 & 15 项检查 & 失败不只在 routing，也在卡片命中 \\
realtime action 通过率 & 33.3\% & 1/3 & 深度模式下仍有 realtime routing 问题 \\
write result count 通过率 & 100\% & 3/3 & FULL\_WRITE 产物数量还在，但内容不稳 \\
memory card 命中率 & 57.1\% & 4/7 & 卡片内容命中明显不稳 \\
ignore rule 命中率 & 100\% & 1/1 & 负样本隔离仍成立 \\
\bottomrule
\end{tabularx}

性能口径：
\begin{itemize}
\item case runtime: 37110.662 ms
\item avg realtime latency: 0.030 ms
\item avg write latency: 12368.586 ms
\item p95 write latency: 15970.542 ms
\end{itemize}

环境与退化信号：
\begin{itemize}
\item 日志出现 \texttt{Stage 2 unavailable, falling back to Jaccard}
\item 日志出现 \texttt{Graphiti not initialized}
\item 本次是 filtered run，所以 case-level deep eval 被跳过
\end{itemize}

\subsection{深度层失败解释}

\begin{tabularx}{\textwidth}{@{}lX@{}}
\toprule
\textbf{batch} & \textbf{失败含义} \\
\midrule
\texttt{batch\_016\_product\_launch\_ops} & 除了 realtime 误判外，公告措辞卡片没有稳定命中“辅助初筛 / 不要写完全替代人工” \\
\texttt{batch\_018\_customer\_policy\_boundary} & 客服口径卡片没有同时覆盖“仅供参考 / 人工复核 / 不要承诺 / 稍后重试” \\
\texttt{batch\_022\_release\_policy\_time\_boundary} & realtime 误判仍在，同时“辅助初筛”口径卡片未稳定命中 \\
\bottomrule
\end{tabularx}

这层能得出的结论：
\begin{itemize}
\item benchmark 已经能把写入侧内容问题打出来，不只是测 routing
\item 当前 FULL\_WRITE 下，发布运营类场景的记忆卡片稳定性不足
\item 但这层\textbf{仍不能}证明 Graphiti / Stage 2 / 完整 case-level long-term recall 的最终水平
\end{itemize}

\section{最终判断}

\textbf{可以肯定的：}
\begin{itemize}
\item 当前 benchmark\_v2 已经具备结构化、分层次、可复跑、能暴露真实问题的基本条件
\item 当前样本不再只是“能跑通”，而是已经能有效区分普通样本和 near-miss / multi-intent 样本
\item 对你们当前阶段最有价值的是：它已经足够拿来压 realtime 分类与写入卡片质量
\end{itemize}

\textbf{不能过度声称的：}
\begin{itemize}
\item 不能说“冲突更新已经验证成功”，因为 relation 命中率这轮是 skipped
\item 不能说“终态 long-term recall 已完整验证”，因为没有完整 case-level deep eval
\item 不能说“系统已生产级可用”，因为 FULL\_WRITE 在非工程场景子集上仍明显失败，且 Graphiti 仍未初始化
\end{itemize}

\textbf{下一步最该做的：}
\begin{itemize}
\item 先用这 3 个轻量失败 batch 回归修 realtime 规则：发布窗口、风险预案、FAQ 截止的边界
\item 再用 \texttt{oc\_launch\_ops} 子集压写入卡片质量，先把 7 个 memory card 检查中的 4/7 提上来
\item Graphiti / Stage 2 初始化稳定后，再补一轮完整 deep case，届时再谈 relation / evidence / final recall 总结论
\end{itemize}

\end{document}
